\section{Broader Impacts}
\label{sec:broader}

\paragraph{Education and mentoring.}
The PIs have a substantial record in education and curriculum development.
PIs Kruegel and Vigna teach several classes related to the proposed research,
and these courses are routinely among the highest-ranked in their
departments. For instance, PI Vigna received the UCSB Academic Senate
Distinguished Teaching Award, the highest teaching award at UC Santa Barbara,
and in 2025 PI Kruegel received the UCSB Outstanding Graduate Mentor Award.
PIs Doup\'e and Bao co-teach ASU's undergraduate ``Introduction to
Cybersecurity'' course, required for all Computer Science students and
averaging over 700 students a semester. The PIs plan to introduce modules
focused explicitly on infrastructure security into their classes, including
lab exercises where students find and exploit defects in complex
infrastructure and propose solutions to the flaws they identify.

PI Doup\'e is a lead developer of the online educational platform
pwn.college and, together with PI Bao, an Associate Director of ASU's
American Cybersecurity Education Institute, which runs the publicly available
pwn.college instance and facilitates classes for over 20 colleges around the
United States~\cite{nelson26:pwncollegefiveyears,shoshitaishvili26:luminarium,nelson25:sensai}.
pwn.college is an open-source system where anyone can learn hands-on
cybersecurity material with only a browser, and the PIs plan to create
open-source and reusable modules for it that focus on infrastructure
security, giving the educational output of this project the same distribution
as the platform itself.

The PIs also have a history of including undergraduate and high school
students in their research, with students often listed as co-authors on
publications. For example, Audrey Dutcher, Paul Grosen, and Jessie Grosen
each joined the lab as high school interns and later continued as
undergraduates at UCSB and ASU. In addition, PIs Kruegel and Vigna have been
faculty mentors in the Early Research Scholars Program, an NSF-supported
research apprenticeship that supports students in their first research
experience, and a letter of collaboration from Ziad Matni, who leads the
program, is attached to this proposal. The PIs also run a summer internship
program that has hosted more than 100 undergraduates over the past decade.

\paragraph{Community outreach: the twinCTF competition.}
The PIs have created the first distributed educational capture-the-flag
competition, called the iCTF. The iCTF has run every year since 2003, and in 2026 it
introduced a new type of challenge that required participants to develop
agents to solve security problems. As part of this project, the PIs
will organize an annual public competition, called twinCTF, centered on
digital twins. In this competition, the participants will be given meta-information bundles
that describe computer infrastructures based on open-source software, and
they will need to develop digital twins that represent those infrastructures.
The organizers, in turn, will run attacks known to work against the original
but localized on the submitted twin, and if an attack succeeds, the
participants get points for having created a representative twin. The
challenges are drawn from the reference networks of
Section~\ref{sec:evaluation} and are released publicly afterward, so that the
competition also grows the benchmark. twinCTF will be open to the public, and
its low barrier to entry and range of difficulty are designed to include
students with different backgrounds and skill levels.

\paragraph{Technology transfer, open source, and dissemination.}
The PIs have a long history of making their research artifacts publicly
available, including Anubis~\cite{anubis}, Wepawet~\cite{cova10:wepawet}, and
angr~\cite{angr,shoshitaishvili2016state}, a binary analysis framework with
more than 9K stars on GitHub. The PIs plan to make \sysname, the
reference-network benchmark, and the twinCTF infrastructure available as open
source under the process described in Section~\ref{sec:buildtest}, and to
contribute the improvements to the anchor products upstream, so that they
reach users who never run \sysname at all. The PIs also have a successful
track record of working with industry partners, and PIs Kruegel and Vigna
founded Lastline, a network security company acquired by VMware in 2020. Our
results will be published in top-tier security venues, such as the IEEE
Symposium on Security and Privacy, the ACM Conference on Computer and
Communications Security, the USENIX Security Symposium, and the Network and
Distributed Systems Security Symposium.
